\chapter{总结与展望}

\section{工作总结}

本文围绕数据选择这一核心方法论，针对深度学习全流程中的三个关键子问题，分别提出了一个框架或方法。

\textbf{CADS：计算预算感知的数据选择（第三章）。}针对训练阶段数据选择与计算资源约束的交互，本文提出CADS框架。CADS首次将计算预算纳入数据选择的双层优化框架，通过概率化重参数化的策略梯度绕过内层非收敛的困难，并通过惩罚松弛与一维损失近似实现高效求解。CADS的核心发现是：\textbf{最优数据子集是计算预算的函数}——低预算偏好简单干净的数据，高预算受益于多样化的数据。在视觉分类和语言模型微调任务上，CADS在所有计算预算下一致领先于基线方法。

\textbf{SWaST：数据选择与权重剪枝的协同优化（第四章）。}针对训练阶段数据选择与模型容量约束的交互，本文提出SWaST框架。SWaST首次透明地分析了数据冗余与参数冗余的双向耦合关系，揭示了独立优化两类冗余的局限性。通过交替执行权重剪枝与核心子集选择，SWaST建立了协同效应：精选数据帮助剪枝（减少过拟合，使权重重要性估计更准确），剪枝后的紧凑模型帮助数据选择（降低特征空间有效维度，使核心子集选择更精准）。SWaST同时识别并解决了"双重损失"这一独特挑战，通过状态保护机制防止不可逆退化。

\textbf{NAP：基于神经激活先验的分布外检测（第五章）。}针对推理阶段数据选择与分布偏移的交互，本文提出NAP方法。NAP发现了一个被现有方法忽视的OOD检测信号：训练使模型在ID数据上形成尖锐的通道内激活先验——少量空间位置产生高激活，其余位置保持低激活。这一先验在OOD样本上不成立（激活更分散），因此可直接用于推理时的OOD判断。NAP设计的即插即用评分函数与现有方法正交互补，在CIFAR和ImageNet基准上与多种基线组合后均实现显著提升。


\section{统一视角}

\subsection{阶段覆盖的完整性}

三项工作覆盖了数据选择在深度学习流程中的三个关键环节：训练阶段的子集选择（CADS）、训练阶段的数据-模型联合选择（SWaST）、推理阶段的数据有效性判断（NAP）。这一覆盖从"选什么数据训练"出发，经过"数据与模型如何协同选择"，到达"如何判断输入是否处于有效域内"，形成了从训练到推理的完整链条。

\subsection{交互视角的统一性}

表~\ref{tab:unified-comparison}从交互维度统一对比三项工作。

% [表6-1: 三项工作的统一对比总结]
\begin{table}[htb]
  \centering
  \caption{三项工作的统一对比——研究问题、阶段定位、交互变量、核心方法与关键发现}
  \label{tab:unified-comparison}
  \fbox{\parbox{0.9\textwidth}{\centering\vspace{2cm}
    \textbf{占位表：三项工作的统一对比总结}\\[0.5em]
    列：工作名称、研究问题、阶段、交互变量、核心方法、关键发现
  \vspace{2cm}}}
\end{table}

三项工作揭示了一个共性的方法论启示：数据选择的最优策略不是孤立的，而是取决于其所处的约束环境。计算资源约束改变了最优训练数据的组成（CADS），模型容量约束与数据冗余相互耦合、需要协同而非独立处理（SWaST），推理时分布偏移要求对数据有效性做在线判断、连接了训练时的选择与推理时的检测（NAP）。

这一共性可以提炼为：\textbf{未来的数据选择研究应系统考虑选择策略与其交互因素的依赖关系，而非将数据选择作为孤立的优化问题。}


\section{局限性}

\subsection{CADS的局限}

CADS的一维损失近似$l(|\mathbf{m}|)$采用的对数线性假设在极端数据分布下可能不够准确，需要更灵活的拟合方法\cite{kaplan2020scaling}。当前实验规模（最大DomainNet约0.5M样本）尚未在十亿级预训练数据上验证，CADS在更大规模上的表现有待探索。此外，CADS在每个计算预算下独立优化选择策略，未利用不同预算间可能存在的共享信息。% [VERIFY: 原论文limitations]

\subsection{SWaST的局限}

状态保护机制引入了额外的前向传播和KL散度计算开销，在极大规模数据集上可能成为瓶颈。当前SWaST仅与RigL\cite{evci2020rigging}剪枝算法验证了兼容性，与其他剪枝方法（如结构化剪枝\cite{he2018soft}）和其他压缩技术（量化\cite{han2016deep, courbariaux2016binarized}、知识蒸馏）的兼容性未经探索。交替频率$\mathcal{R}$和保护强度$\lambda$的调优目前依赖经验搜索，缺乏理论指导。% [VERIFY: 原论文limitations]

\subsection{NAP的局限}

NAP依赖全局池化层（CNN）或注意力机制（Transformer）的存在，对不含此类结构的架构不直接适用。在近分布OOD（near-OOD）场景下——即OOD数据与ID数据在语义上高度相似时——通道内激活模式的差异也会缩小，检测区分度有限。组合权重$w$需要通过伪OOD数据调优，不同ID/OOD对的最优$w$可能不同。% [VERIFY: 原论文limitations]


\section{未来研究方向}

\subsection{跨阶段的数据选择联合建模}

本文三项工作分别处理不同阶段的数据选择，但它们之间存在潜在联系：训练阶段的数据选择决定了模型的有效域，这直接影响推理时OOD检测的难度。一个有意义的未来方向是：能否在训练阶段做数据选择时，直接优化推理时的OOD检测能力？例如，选择能使模型学到更清晰决策边界的训练子集，从而在推理时更容易区分ID与OOD输入。更广义地，端到端考虑从数据选择到模型训练到推理部署的全链条优化，是一个值得探索的方向。

\subsection{大模型时代的数据选择新挑战}

LLM预训练面临万亿级数据的选择与配比问题\cite{albalak2024survey}。CADS的预算感知思路和来源级选择（CADS-S）可为此提供新视角：在不同训练阶段（预训练初期 vs. 后期），最优数据配比是否也应随之变化？LLM微调中的数据选择\cite{zhou2023lima}还需要在数据质量、多样性和安全性（alignment）之间寻求平衡，计算预算约束更为严格。此外，多模态大模型中跨模态数据的配比和互补性选择，也是数据选择领域面临的新挑战。

\subsection{更多交互维度的探索}

本文探索了数据选择与三个交互变量（计算资源、模型容量、推理时分布偏移）的关系。类似的交互可能存在于更多维度：数据选择与模型架构搜索（NAS\cite{liu2018darts}）之间是否存在耦合？数据选择与其他压缩技术（量化\cite{han2016deep}、知识蒸馏）之间是否存在类似SWaST发现的协同关系？在持续学习场景中，数据分布随时间变化\cite{borsos2020coresets}，数据选择策略需要如何适应？在联邦学习场景中，数据分布于多个参与方，又该如何协调选择？这些问题构成了数据选择领域的广阔未来图景。
